\documentclass[11pt]{amsart}
\usepackage[T1]{fontenc}
\usepackage[utf8]{inputenc}
\usepackage{amssymb}
\usepackage[margin=1.05in]{geometry}
\usepackage[hidelinks]{hyperref}
\newcommand{\stamptime}{\the\numexpr\time/60\relax:\ifnum\numexpr\time-60*(\time/60)\relax<10 0\fi\the\numexpr\time-60*(\time/60)\relax}
\AddToHook{shipout/foreground}{\put(0.5\paperwidth,-\paperheight+0.55cm){\makebox[0pt]{\tiny\ttfamily unit/doc/\jobname.tex \ --- \ compiled \today, \stamptime}}}
\newtheorem{theorem}{Theorem}
\newtheorem{proposition}[theorem]{Proposition}
\newtheorem{lemma}[theorem]{Lemma}
\newtheorem{corollary}[theorem]{Corollary}
\theoremstyle{definition}
\newtheorem{definition}[theorem]{Definition}
\theoremstyle{remark}
\newtheorem{remark}[theorem]{Remark}
\DeclareMathOperator{\Syz}{Syz}
\DeclareMathOperator{\res}{res}
\newcommand{\sm}{\mathrm{span}}
\title{Constructive Artin--Rees, and the ACL certificate from it}
\date{28 October 2026}
\begin{document}
\maketitle

\noindent\textbf{Status.}
This note states the constructive Artin--Rees theorem as it is checked by
Unit in \texttt{unit/sag/artin-rees}, and explains how the ZMT certificate
of ACL [ACL, \S4] follows from it. Proofs of the theorem itself are in
\texttt{artin-rees.tex} (\S\S1--5); numbers ``AR~$n$'' refer to that note.
Every statement below is checked by Unit, except where marked
\emph{(paper)}.

\section{The constructive Artin--Rees theorem}

Rings are commutative, unital and sets. ``Merely'' means: under
propositional truncation.

\begin{definition}[AR~1]
A ring $R$ is \emph{CFN} if
\begin{itemize}
\item[(D)] for $u_1,\dots,u_k,v\in R$ it is decidable whether $v\in\sm(u)$;
\item[(C)] for $u_1,\dots,u_k\in R$ the module $\Syz(u)$ is merely
  generated by finitely many syzygies;
\item[(N)] $R$ is merely the image of a presentation by codes that is
  noetherian in the sense of Coquand--Persson [CP].
\end{itemize}
\end{definition}

CFN is a proposition. Discrete fields, $\mathbb Z$, $R[X]$, quotients by
finitely generated ideals and localisations $R[1/f]$ of CFN rings are CFN
(AR~10, 12, 13). For example $A=k[b][a]$, with $k$ a discrete field, is
CFN. The class is therefore large enough for finitely presented algebras
over $\mathbb Z$ or over a discrete field. No choice principle and no
excluded middle are used.

Powers of $I=(a,b)$ acting on a submodule $P$ are defined inductively:
$I^0P=P$, and $x\in I^{m+1}P$ if $x=ax_1+bx_2$ with $x_1,x_2\in I^mP$.
An element of $I^mP$ is thus a binary tree of depth $m$ whose leaves lie
in $P$.

\begin{theorem}[Artin--Rees, AR~15]\label{thm:ar}
Let $R$ be CFN, $a,b\in R$, $I=(a,b)$, and
$N=\sm(n_1,\dots,n_s)\subseteq R^r$. There merely is $c\in\mathbb N$ such
that for every $m$ and every $x\in R^r$
\[
 x\in I^{m+c}R^r\cap N\ \Longrightarrow\ x\in I^m\bigl(I^cR^r\cap N\bigr).
\]
\end{theorem}

The converse inclusion is immediate. The quantifiers come in the order
$\exists c\,\forall m\,\forall x$: one bound serves for every $m$. This
uniformity is what the application needs, since $m$ is chosen after $c$.

\emph{Idea of the proof.} In the bigraded ring $R'=R[U][V]$ (CFN by AR~10
applied twice), consider the module $Q^+\subseteq R'^r$ generated by the
$n_i$ and by $(U-aY)e_k$, $(V-bY)e_k$. By the degree corollary (AR~9),
$Q^+\cap R'^r$ is merely generated by finitely many $G_1,\dots,G_p$. The
bound $c$ is the sum of their total degrees. A ``good'' element splits in
degree $m+c+1$ into two halves of degree $m+c$, and this gives the step
$I^{m+c+1}\cap N\subseteq I(I^{m+c}\cap N)$. The bound comes from
noetherianity through (N); (D) and (C) make the generators $G_i$
computable.

\begin{remark}
In Unit this is \texttt{ArtinRees.Over.artin-rees :
Trunc[Bound]}, where \texttt{Bound} is the record of $c$ and the
implication above. The statement for a finitely presented $M$ and a
finitely presented submodule reduces to this one
(\emph{paper}, AR~\S5: write $M=R^r/K$ and take the preimage of $N$,
which contains $K$). The principal case $I=(f)$, $M=R$, used for flatness rather than ACL, is
proved separately (AR~14, \texttt{ArtinReesPrincipal}), for the uniform
Kill lemma (AR~27).
\end{remark}

\begin{corollary}[coefficients, AR~26]\label{cor:coef}
Under the hypotheses of Theorem~\ref{thm:ar}, if
$x\in I^{m+c}R^r\cap N$, then merely $x=\sum_j\lambda_jn_j$ with every
$\lambda_j\in I^m$.
\end{corollary}

\begin{proof}
The theorem gives a tree of depth $m$ whose leaves lie in $N$. Each leaf
is a combination of the $n_j$, and at a node $ax_1+bx_2$ the
coefficients combine to $a\lambda_{1,j}+b\lambda_{2,j}\in I^{m+1}$.
(\texttt{PowSpan.Free.coefficients}.)
\end{proof}

This is the form in which Artin--Rees is used below. A combination
that lies \emph{deep} in $I$ can be rewritten with coefficients that lie
deep in $I$.

\section{The ACL statement}

Let $A$ be a gcd domain (in the sense of \texttt{sag/gcd}) that is CFN,
for instance $A=k[b][a]$. Let $a,b\in A$, $I=(a,b)$, and
\[
 p=x+af+bg,\qquad q=y+au+bv\qquad(f,g,u,v\in A[x,y]).
\]
Let $P,Q\in A[X,Y,Z]$ be their homogenisations, of degrees $d,e$. Let
$L$ be any $A$-algebra with elements $x,y$ satisfying $p(x,y)=q(x,y)=0$.

\begin{theorem}[ZMT certificate, ACL \S4]\label{thm:acl}
There merely are $s,r,t\in L$ with
\[
 s=1+ar+bt,
\]
such that $s$, $sx$ and $sy$ are each roots of a monic polynomial over
$A$.
\end{theorem}

Geometrically: the fibre of $\operatorname{Spec}L\to\operatorname{Spec}A$
over $a=b=0$ is the single point $x=y=0$, and the theorem gives an
element $s\equiv1$ modulo $I$ such that $L[1/s]$ is finite over $A$
there. This is the constructive content of Zariski's Main Theorem
in this situation. In Unit the statement is \texttt{Certificate} of
\texttt{ACLZMTCong}, and the proof for a CFN gcd domain is
\texttt{ACLCertificateGeneral.certificate}. Its instance for $k[b][a]$ is
\texttt{ACLCertificateK}.

\section{From Artin--Rees to the certificate}

The proof works in an explicit Čech model of the twisted structure
sheaves of the projective closure $W=V(P,Q)\subseteq\mathbb P^2_A$.
Here $K^j(n)$ is the module of $j$-cochains of degree $n$ for the
three standard charts, with coefficients Laurent polynomials, and
\[
 N^j(n)=P\,K^j(n-d)+Q\,K^j(n-e)
\]
is the module of cochains divisible by the equations.
Let $\res$ be restriction to the chart $Z$, followed by evaluation at
$(x,y)$; it sends $N^0$ to $0$. There is a fixed $0$-cochain
$z^*\in K^0(-1)$ (the cochain $Z^{-1}$) with $\res(Zz^*)=1$, but $z^*$ is
\emph{not} closed modulo $N^1$. The proof has five steps; only step~3
uses Artin--Rees.

\medskip\noindent\textbf{Step 1: approximations of $z^*$ (data (F1)).}
Iterating $Z^{-1}=E\,Z^{-1}+P\,u$, with $E$ a combination of $a$ and $b$,
gives the following for every $n\ge1$: vectors $\rho_n\in A^r$ with
coordinates in $I^{n-1}$, and $c_n\in I^nK^0(-1)$, with
\[
 \delta_0z^*\equiv G\rho_n+\delta_0c_n\pmod{N^1}.\tag{$\beta_n$}
\]
Here $G\rho=\sum_t\rho_tm_t$ is a combination of fixed generators $m_t$
of the closed $1$-cochains modulo coboundaries (\texttt{CechF1},
\texttt{FixedCochain}). In words, the obstruction to closing $z^*$ is
$G\rho_n$, a combination with coefficients in $I^{n-1}$.

\medskip\noindent\textbf{Step 2: the module of liftable obstructions is
finitely generated.} Let
\[
 N=\{\rho\in A^r\mid \exists v\in K^0(-1),\ G\rho-\delta_0v\in N^1\}.
\]
By $(\beta_n)$, $\rho_n\in N$, with $v=z^*-c_n$. The gcd data of $P,Q$
(the common factor $D_0$, cofactors $P',Q'$, and the conditions (S0) and
(S1) of \texttt{HomLambda.Data}) describe $N$ as a projection of the
kernel of a finite matrix $\mathcal M$ over $A$ (AR~29, AR~30). Coherence of $A$ then gives merely
$u_1,\dots,u_p$ and $b_1,\dots,b_p\in K^0(-1)$ with
\[
 N=\sm(u_1,\dots,u_p),\qquad Gu_j-\delta_0b_j\in N^1
\]
(\texttt{ACLRelations.corN}, AR~31).

\medskip\noindent\textbf{Step 3: the correction, by Artin--Rees.}
Let $c$ be an Artin--Rees bound for $N\subseteq A^r$
(Theorem~\ref{thm:ar}). Fix $m$ and let $n=c+m+1$. Then
$\rho_n\in I^{m+c}A^r\cap N$, so Corollary~\ref{cor:coef} gives
$\rho_n=\sum_j\lambda_ju_j$ with $\lambda_j\in I^m$. Put
\[
 w=c_n+\sum_j\lambda_jb_j\in I^mK^0(-1),\qquad h=z^*-w.
\]
Modulo $N^1$,
$\delta_0w\equiv\delta_0c_n+\sum_j\lambda_jGu_j=\delta_0c_n+G\rho_n\equiv\delta_0z^*$.
So
\[
 \delta_0h\in N^1,\qquad h-z^*\in I^mK^0(-1)
\]
(\texttt{ACLCorrection.correction}). With $m=1$:
$h=z^*+am_1+bm_2$ is closed modulo $N^1$.

This is where the uniformity of the bound matters. Step~2 alone gives
$\rho_n=\sum_j\mu_ju_j$ with arbitrary coefficients $\mu_j$; then
$\sum_j\mu_jb_j$ corrects $\delta_0z^*$, but the correction is no longer
small in $I$. Artin--Rees turns ``$\rho_n$ is deep in $I$ and lies in
$N$'' into ``$\rho_n$ is a combination of generators of $N$ with
coefficients deep in $I$''. That is exactly what keeps $w$ in $I^m$.

\medskip\noindent\textbf{Step 4: the element $s$.} Multiplication by
$X,Y,Z$ sends $h$ to $0$-cochains $Xh,Yh,Zh$ of degree~$0$, which are
still closed modulo $N^1$. Restriction gives
\[
 s=\res(Zh)=\res(Zz^*)+a\res(Zm_1)+b\res(Zm_2)=1+ar+bt,
\]
\[
 \res(Xh)=sx,\qquad \res(Yh)=sy
\]
(\texttt{ACLCertificate.From}).

\medskip\noindent\textbf{Step 5: integrality.} Closed $0$-cochains of
degree $0$ modulo $N^0$ form a ring, and modulo coboundaries they form a
finitely generated $A$-module: data (F2) give coordinates
$\kappa:\ \text{closed sections}\to A^{s'}$ (\texttt{CechSections}). For
a closed section $\gamma$, apply the dependence lemma (AR~32:
in a module with (D) and the principle of free lists, a sequence $f_n$
merely has some $f_N\in\sm(f_0,\dots,f_{N-1})$) to $f_n=\kappa(\gamma^n)$
in $A^{s'}$, which satisfies both hypotheses because $A$ is CFN (AR~4).
This gives $\gamma^N-\sum_{i<N}x_i\gamma^i\in N^0$. Restricting,
$\res(\gamma)$ is a root of a monic polynomial over $A$
(\texttt{SeqDependence}, \texttt{ACLCertificate.Walk}). Applied to
$\gamma=Zh,Xh,Yh$ this gives Theorem~\ref{thm:acl}.

\medskip
Each step produces a proposition from merely given data (bound,
generators, coefficients, dependence). These are eliminated one after
the other into the truncated conclusion, so no choice is needed.

\section{Where noetherianity enters}

Noetherianity (N) of $A$ enters twice, both times through CFN:
\begin{itemize}
\item in step~3, through Artin--Rees (Theorem~\ref{thm:ar}: degree
  corollary for $R'=A[U][V]$);
\item in step~5, through the principle of free lists on $A^{s'}$.
\end{itemize}
Before this work, both steps used a coded noetherian presentation of
$A$ itself (\texttt{ReesLifting}, \texttt{GradedDependence}). The
certificate now depends only on ``$A$ is a CFN gcd domain''.

\begin{thebibliography}{9}
\bibitem{ACL} M.~E.~Alonso, T.~Coquand, H.~Lombardi, Revisiting Zariski
  main theorem from a constructive point of view, \emph{J.~Algebra} 406
  (2014), 46--68.
\bibitem{CP} T.~Coquand, H.~Persson, Gröbner bases in type theory,
  TYPES'98, LNCS 1657 (1999).
\bibitem{R} W.~Ruitenburg, Constructions of finitely generated submodules
  of constructively noetherian modules, \emph{Compositio Math.} 62 (1987),
  47--52.
\end{thebibliography}
\end{document}
